\begin{figure*}[t]
\centering
\includegraphics[width=\textwidth]{figures/overview_vector.pdf}
\caption{Overview: kernel equivalence, dimension, and execution cost.
(a) Schematic ideal-expectation geometry. A common input-independent
state transformation preserves the specified kernel when it induces
the displayed $M$-isometry of features; arbitrary unitaries need not
do so. $K_\theta$ denotes the sampled kernel matrix, distinct from
the coefficient Gram matrix $G$ in (b).
(b) For a fixed data basis with $\{1,h_1,\ldots,h_A\}$ linearly
independent on the actual input domain, fixed $M\succ0$, fixed
$\gamma>0$, and analytic coefficients on an open connected parameter
box, Theorem~\ref{thm:gram} characterizes all-input equivalence,
global coordinate inactivity, and generic differential rank through $G$.
The two-qubit local-XYZ Kyriienko example has six source parameters
and rank four, certified by a Gram minor and an output gauge.
(c) The ChebyshevRx example separates retained symbols ($7\to5$),
exact smooth kernel dimension ($r=4$), and globally normalized
angle tasks ($3\to3$ in the boundary and transport arms;
Table~\ref{tab:global}). Its parameter--data products lie outside
(b)'s fixed-basis construction: the separate closed-form factorization
and exact minor establish its dimension. Task counts are
schedule-specific, not native-gate or hardware-shot measurements.}
\label{fig:overview}
\end{figure*}
